% Source: Wald GR, physical PDF pp. 246--250
%         (printed pp. 233--237).
% Coverage: Section 9.4, Existence of Maximum Length Curves; equations (9.4.1)--(9.4.9).
% Figures/tables/footnotes: Figure 9.5; no tables or footnotes.
% Status: source-reviewed and compile-clean.
% Uncertainties: none.

\subsection{Existence of Maximum Length Curves}\label{sec:9.4}

\setcounter{theorem}{0}

In the previous two sections we established by means of ``local'' calculations
necessary criteria for a timelike curve to be a curve of maximum length between two
points (Theorem~9.3.3) or between a point and a hypersurface \(\Sigma\)
(Theorem~9.3.5) as well as conditions under which these criteria could not be met
(Propositions~9.3.1 and 9.3.4). In this section, we shall use global arguments
involving compactness of the spaces of causal curves \(\mathcal C(p,q)\) and
\(\mathcal C(\Sigma,q)\) defined in Section~8.3 to prove existence of maximum length
curves in globally hyperbolic spacetimes. In the next section we shall see that,
under further hypotheses, these two sets of results lead to contradictions if all
geodesics are complete, thereby producing the singularity theorems.

The length, \(\tau\), of a smooth (or even \(C^1\)) causal curve \(\lambda\) between
points \(p,q\in M\) with tangent \(T^a=(\partial/\partial t)^a\) is given by the formula

\begin{equation}
  \tau[\lambda]=\int(-T^aT_a)^{1/2}\,\dd t .
  \tag{9.4.1}\label{eq:9.4.1}
\end{equation}

However, for the global arguments of this section, it is necessary to generalize this
notion of length to continuous causal curves between \(p\) and \(q\) in order that
\(\tau\) be defined for all curves in \(\mathcal C(p,q)\). Let
\(\widetilde{\mathcal C}(p,q)\) denote the subset of \(\mathcal C(p,q)\) consisting of
smooth timelike curves, with the topology induced by \(\mathcal C(p,q)\). Then---with
the possible exception of cases where null geodesics connect \(p\) and
\(q\)---\(\widetilde{\mathcal C}(p,q)\) is dense in \(\mathcal C(p,q)\), i.e., except
for certain null geodesics every continuous causal curve can be expressed as a limit
[in the topology of \(\mathcal C(p,q)\)] of a sequence of smooth timelike curves. If
\(\tau\) were continuous on \(\widetilde{\mathcal C}(p,q)\), we could extend it to a
continuous function on all of \(\mathcal C(p,q)\) by setting
\[
  \tau[\mu]=\lim_{n\to\infty}\tau[\lambda_n] ,
\]
where \(\{\lambda_n\}\) is a sequence in \(\widetilde{\mathcal C}(p,q)\) which
approaches the continuous causal curve \(\mu\in\mathcal C(p,q)\). However, as
Figure~\ref{fig:9.5} illustrates, \(\tau\) is not continuous on
\(\widetilde{\mathcal C}(p,q)\). Arbitrarily close in the topology of
\(\mathcal C(p,q)\) to any smooth timelike curve we can find a ``zigzag'' smooth
timelike curve of length arbitrarily close to zero. Nevertheless, we shall show below
that \(\tau\) is \emph{upper semicontinuous} on \(\widetilde{\mathcal C}(p,q)\), i.e.,
for each \(\lambda\in\widetilde{\mathcal C}(p,q)\) given \(\epsilon>0\) there exists an
open neighborhood \(O\subseteq\widetilde{\mathcal C}(p,q)\) of \(\lambda\) such that for
all \(\lambda'\in O\) we have \(\tau[\lambda']\leq\tau[\lambda]+\epsilon\). Given that
\(\tau\) is upper semicontinuous on \(\widetilde{\mathcal C}(p,q)\), we can extend it
to an upper semicontinuous function on \(\mathcal C(p,q)\) as follows. For
\(\mu\in\mathcal C(p,q)\) and \(O\subseteq\mathcal C(p,q)\) an open neighborhood of
\(\mu\), we define

\begin{equation}
  \tau[O]=\operatorname{l.u.b.}\{\tau[\lambda]\mid
    \lambda\in O,\ \lambda\in\widetilde{\mathcal C}(p,q)\} ,
  \tag{9.4.2}\label{eq:9.4.2}
\end{equation}

where ``l.u.b.'' denotes the least upper bound. Then we define \(\tau[\mu]\) by,

\begin{equation}
  \tau[\mu]=\operatorname{g.l.b.}\{\tau[O]\mid
    O\text{ an open neighborhood of }\mu\} ,
  \tag{9.4.3}\label{eq:9.4.3}
\end{equation}

where ``g.l.b.'' denotes the greatest lower bound. Thus, the key fact that allows us
to extend the definition of \(\tau\) to \(\mathcal C(p,q)\) is expressed by the following
proposition.

% Source: Wald GR, physical PDF p. 247 (printed p. 234), Figure 9.5.
% Coverage: A smooth timelike curve and nearby zigzag curve.
% Figures/tables/footnotes: Figure 9.5.
% Status: source-reviewed and compile-clean.
% Uncertainties: none.

\begingroup
\renewcommand{\thefigure}{9.5}
\begin{figure}[tb]
  \centering
  \begin{tikzpicture}[x=1cm,y=1cm,>=Stealth,thick]
    \coordinate (p) at (0,0);
    \coordinate (q) at (.55,5.2);
    \draw (p) .. controls (-.02,1.7) and (.03,3.65) .. (q);
    \draw (-.04,.18)
      .. controls (.9,.58) and (-.75,.97) .. (-.04,1.43)
      .. controls (.75,1.88) and (-.62,2.25) .. (-.03,2.73)
      .. controls (.68,3.18) and (-.45,3.56) .. (.12,4.01)
      .. controls (.66,4.42) and (.2,4.82) .. (.55,5.06);
    \fill (p) circle (2pt) node[below] {\(p\)};
    \fill (q) circle (2pt) node[above] {\(q\)};
    \draw[->] (-1.25,4.22) to[out=8,in=174] (.29,4.18);
    \node[left] at (-1.22,4.22) {\(\mu\)};
    \draw[->] (1.72,4.55) to[out=190,in=18] (.55,4.43);
    \node[right] at (1.72,4.55) {\(\mu'\)};
  \end{tikzpicture}
  \caption{A smooth timelike curve \(\mu\) from \(p\) to \(q\). Arbitrarily close
  to \(\mu\) in the topology of \(\mathcal C(p,q)\) is a smooth timelike curve
  \(\mu'\) with total elapsed proper time arbitrarily close to zero.}
  \label{fig:9.5}
\end{figure}
\endgroup


\begin{proposition}\label{prop:9.4.1}
Let \((M,g_{ab})\) be a strongly causal spacetime. Let \(p,q\in M\) with
\(q\in I^+(p)\). Then \(\tau\) is upper semicontinuous on
\(\widetilde{\mathcal C}(p,q)\).
\end{proposition}

\begin{proof}
Let \(\lambda\in\widetilde{\mathcal C}(p,q)\). We parameterize \(\lambda\) by proper
time and denote its tangent by \(u^a\). Within a normal neighborhood of each point
\(r\in\lambda\), the spacelike geodesics orthogonal to \(u^a\) will form a
three-dimensional spacelike hypersurface. Within a sufficiently small open
neighborhood \(U\subseteq M\) of \(\lambda\), these hypersurfaces will foliate \(U\);
i.e., a unique hypersurface will pass through each point of \(U\). On \(U\), we define
the function \(F\) by setting \(F(p)\) equal to the proper time value of \(\lambda\) at
the intersection of \(\lambda\) with the hypersurface on which \(p\) lies. Then
\(\nabla^aF\) is timelike everywhere in \(U\), and on \(\lambda\) we have
\(u^a=-\nabla^aF\), so \(\nabla^aF\nabla_aF=-1\) on \(\lambda\). Now, let
\(\rho\in\widetilde{\mathcal C}(p,q)\) with \(\rho\subseteq U\). We parameterize
\(\rho\) by \(F\) and denote its tangent by \(v^a\). Thus, by our parameterization
choice, we have

\begin{equation}
  v^a\nabla_aF=1 .
  \tag{9.4.4}\label{eq:9.4.4}
\end{equation}

We decompose \(v^a\) as

\begin{equation}
  v^a=\alpha\nabla^aF+m^a ,
  \tag{9.4.5}\label{eq:9.4.5}
\end{equation}

where \(m^a\nabla_aF=0\), and thus \(m^a\) is spacelike. Contracting
Eq.~\eqref{eq:9.4.5} with \(\nabla_aF\) and using Eq.~\eqref{eq:9.4.4}, we evaluate
\(\alpha\) and find

\begin{equation}
  v^a=\frac{\nabla^aF}{\nabla_bF\nabla^bF}+m^a .
  \tag{9.4.6}\label{eq:9.4.6}
\end{equation}

Therefore, we find

\begin{equation}
  v^av_a=(\nabla^aF\nabla_aF)^{-1}+m^am_a ,
  \tag{9.4.7}\label{eq:9.4.7}
\end{equation}

and thus

\begin{equation}
  (-v_av^a)^{1/2}\leq(-\nabla_aF\nabla^aF)^{-1/2}
  \tag{9.4.8}\label{eq:9.4.8}
\end{equation}

since \(m^a\) is spacelike. Since \(\nabla^aF\) is continuous, given
\(\epsilon>0\), we can choose a neighborhood \(U'\subseteq U\) of \(\lambda\) so that
\((-\nabla^aF\nabla_aF)^{-1/2}\leq1+\epsilon/\tau[\lambda]\) in \(U'\). Then for all
\(\rho\in\widetilde{\mathcal C}(p,q)\) contained in the open neighborhood \(O'\) in
\(\widetilde{\mathcal C}(p,q)\) defined by \(U'\), we have

\begin{equation}
  \tau[\rho]=\int(-v^av_a)^{1/2}\,\dd F\leq\tau[\lambda]+\epsilon ,
  \tag{9.4.9}\label{eq:9.4.9}
\end{equation}

which proves upper semicontinuity.
\end{proof}

In Section~8.3 we defined \(\mathcal C(\Sigma,p)\) for a Cauchy surface \(\Sigma\) in a
globally hyperbolic spacetime, but, more generally, we may define
\(\mathcal C(\Sigma,p)\) analogously for any achronal set \(\Sigma\) in a strongly
causal spacetime. Essentially the same arguments as given above establish that
\(\tau\) is an upper semicontinuous function on the space
\(\widetilde{\mathcal C}(\Sigma,p)\) of smooth timelike curves from \(\Sigma\) to \(p\).
Therefore, by the same argument as given above, one can extend \(\tau\) to an upper
semicontinuous function defined on all of \(\mathcal C(\Sigma,p)\).

In the previous section, we showed that the necessary and sufficient condition for a
\emph{smooth} curve to locally maximize the length between two points or a point and a
hypersurface was that it be a geodesic without conjugate points. Now that we have
extended the definition of \(\tau\) to continuous curves, there is a possibility that a
continuous, nonsmooth curve between two points or a point and a hypersurface could
have length greater than or equal to that of any geodesic. However, this possibility
can be ruled out as follows. By direct calculation, one can prove that in any convex
normal neighborhood \(U\), the unique geodesic \(\gamma\) connecting two causally
related points \(r,s\in U\) has length strictly greater than that of any other piecewise
smooth causal curve connecting those points (see Proposition~4.5.3 of Hawking and
Ellis 1973 for a proof). Therefore, by upper semicontinuity any continuous causal
curve, \(\mu\), connecting \(r\) and \(s\) in \(U\) must satisfy
\(\tau[\mu]\leq\tau[\gamma]\). However, if equality held with \(\mu\ne\gamma\), let
point \(q\) be such that \(q\in\mu\) but \(q\notin\gamma\). Let \(\gamma_1,\gamma_2\)
be the geodesic segments connecting \(r\) and \(q\) and \(q\) and \(s\), respectively.
Since by the above result \(\gamma_1\) maximizes the length between \(r\) and \(q\),
while \(\gamma_2\) maximizes the length between \(q\) and \(s\), we have
\(\tau[\gamma_1]+\tau[\gamma_2]\geq\tau[\mu]=\tau[\gamma]\), which contradicts the
fact that \(\gamma\) has strictly greater length than any other piecewise smooth curve
between \(r\) and \(s\). Thus, within any convex normal neighborhood, the unique
geodesic connecting any pair of causally related points has length strictly greater
than that of any other continuous causal curve connecting the points. Thus, an
arbitrary continuous causal curve connecting any two points cannot be a curve of
maximum length between those points unless it is a geodesic, since if it failed to be
a geodesic at any point, we could deform it in a convex normal neighborhood of that
point to obtain a longer curve. Thus, in view of Theorem~9.3.3 we have the following
result.

\begin{theorem}\label{thm:9.4.2}
Let \((M,g_{ab})\) be a strongly causal spacetime. Let \(p,q\in M\) with
\(q\in I^+(p)\), and consider the length function \(\tau\) defined on
\(\mathcal C(p,q)\). A necessary condition for \(\tau\) to attain its maximum value at
\(\gamma\in\mathcal C(p,q)\) is that \(\gamma\) be a geodesic with no point conjugate
to \(p\) between \(p\) and \(q\).
\end{theorem}

Similarly, we have:

\begin{theorem}\label{thm:9.4.3}
Let \((M,g_{ab})\) be a strongly causal spacetime. Let \(p\in M\), let \(\Sigma\) be
an achronal, smooth spacelike hypersurface, and consider the length function \(\tau\)
defined on \(\mathcal C(\Sigma,p)\). A necessary condition for \(\tau\) to attain its
maximum value at \(\gamma\in\mathcal C(\Sigma,p)\) is that \(\gamma\) be a geodesic
orthogonal to \(\Sigma\) with no point conjugate to \(\Sigma\) between \(\Sigma\) and
\(p\).
\end{theorem}

The above results, of course, do not assert that \(\tau\) must attain a maximum
value. However, we conclude this section by proving two key results which show that a
maximum is always attained in globally hyperbolic spacetimes.

\begin{theorem}\label{thm:9.4.4}
Let \((M,g_{ab})\) be a globally hyperbolic spacetime. Let \(p,q\in M\) with
\(q\in J^+(p)\). Then there exists a curve \(\gamma\in\mathcal C(p,q)\) for which
\(\tau\) attains its maximum value on \(\mathcal C(p,q)\).
\end{theorem}

\begin{proof}
By Theorem~8.3.9, \(\mathcal C(p,q)\) is compact, and by
Proposition~\ref{prop:9.4.1}, \(\tau\) is upper semicontinuous. Hence, by a slight
generalization of Theorem~A.6 of Appendix~A, \(\tau\) is bounded and attains its
maximum.
\end{proof}

\begin{theorem}\label{thm:9.4.5}
Let \((M,g_{ab})\) be a globally hyperbolic spacetime. Let \(p\in M\) and let
\(\Sigma\) be a Cauchy surface. Then there exists a curve
\(\gamma\in\mathcal C(\Sigma,p)\) for which \(\tau\) attains its maximum value on
\(\mathcal C(\Sigma,p)\).
\end{theorem}

\begin{proof}
Again, \(\mathcal C(\Sigma,p)\) is compact and \(\tau\) is upper semicontinuous.
\end{proof}
